% figures/fig2-ir-expr.tex — IR expressiveness comparison (Task T11)
\begin{figure}[t]
\centering
\begin{tikzpicture}[
  >=Stealth,
  scale=0.72, every node/.style={scale=0.72},
  panel/.style={draw, rounded corners=2pt, inner sep=4pt, align=left,
                minimum width=3.6cm},
  lbl/.style={font=\footnotesize\bfseries, anchor=south},
  code/.style={font=\scriptsize\ttfamily, align=left, anchor=north west},
  note/.style={font=\scriptsize, align=center},
]
  % --- Left panel: LLVM IR ---
  \node[panel, fill=black!8, minimum height=2.6cm] (left) {};
  \node[lbl] at (left.north) {通用 IR (LLVM)};

  \node[code] at ([shift={(2pt,-3pt)}]left.north west) {%
    \%r = call \{i256, i1\}\\
    \quad @llvm.uadd.with.overflow\\
    \quad\quad .i256(\%x, \%y)\\[1pt]
    \%sum\enspace = extractvalue \%r, 0\\
    \%cout = extractvalue \%r, 1%
  };

  \node[note, anchor=north] at ([yshift=-1pt]left.south) {%
    进位位以多返回值结构间接表达};

  % --- Right panel: dMIR ---
  \node[panel, fill=black!3, minimum height=2.6cm,
        right=0.4cm of left] (right) {};
  \node[lbl] at (right.north) {dMIR};

  \node[code] at ([shift={(2pt,-3pt)}]right.north west) (rcode) {%
    (sum, cout) = ADC(x, y, cin)%
  };

  \node[font=\scriptsize, align=left, anchor=north west]
    at ([yshift=-2pt]rcode.south west) {%
    $: \mathrm{BV}_{256}\!\times\!\mathrm{BV}_{256}\!\times\!\mathrm{BV}_{1}$\\
    $\;\;\to \mathrm{BV}_{256}\!\times\!\mathrm{BV}_{1}$%
  };

  \node[note, anchor=north] at ([yshift=-1pt]right.south) {%
    进位位作为一等位向量算子};

  % --- Contrast arrow ---
  \draw[<->, thick, densely dashed]
    ([yshift=-0.2cm]left.east) -- ([yshift=-0.2cm]right.west)
    node[midway, above, font=\scriptsize] {对比};

\end{tikzpicture}
\caption{通用 IR 与 dMIR 在带进位加法上的表达力对比：左侧 LLVM IR 的 \texttt{uadd.with.overflow} 返回多值结构，右侧 dMIR 把进位作为一等位向量算子。}
\label{fig:ir-expr}
\end{figure}
